\chapter{Méthodologie d'analyse}\label{cap:methodologie}
%--------------------------------------

\section{Démarche documentaire}

J'ai construit la revue à partir des articles sur les agents LLM et les modèles factoriels, des textes officiels de gouvernance et des dépôts GitHub. Les recherches portaient notamment sur les agents financiers, les frameworks quantitatifs, le backtesting et les workflows multi-agents. Les solutions retenues fournissaient une brique d'orchestration, d'analyse financière, d'infrastructure quantitative ou de contrôle, avec du code accessible. Les projets purement promotionnels ou trop éloignés du cas ETF ont été écartés.

Le corpus comprend 17 fiches : 6 frameworks généralistes, 6 briques quantitatives et 5 projets financiers. Huit dépôts ont été retenus pour la due diligence finale. Ce décompte décrit les fiches étudiées, sans prétendre à une revue exhaustive. Les documents mobilisés sont archivés dans \texttt{papers} lorsqu'une version légalement téléchargeable est disponible ; une version de travail des auteurs est privilégiée si le texte de l'éditeur est fermé.

La revue documentaire sert aux choix d'architecture ; les mesures empiriques portent sur les cours quotidiens des quatre ETF et de SPY sur 2015--2025, puis sur les sorties de synthèse et d'allocation. Les fonctionnalités, licences et conditions de maintenance des dépôts restent datées et doivent être revérifiées avant adoption.

\section{Revue réglementaire arrêtée au 30 août 2026}

Cette revue est documentaire et ne constitue pas un avis juridique. Elle sert à relier l'architecture proposée aux sources officielles disponibles à la date d'arrêté.

\begin{itemize}
    \item \textbf{MiFID II, article 17 :} le trading algorithmique implique des systèmes et contrôles de risque adaptés, des limites, des tests et une capacité à décrire les paramètres et contrôles à l'autorité compétente \citep{mifidArticle17}.
    \item \textbf{DORA :} le règlement (UE) 2022/2554 s'applique depuis le 17 janvier 2025 et renforce la résilience opérationnelle numérique des entités financières \citep{dora2022}.
    \item \textbf{AI Act :} le règlement (UE) 2024/1689 s'applique en principe depuis le 2 août 2026, avec des dates spécifiques pour certaines dispositions, notamment le 2 février 2025, le 2 août 2025 et le 2 août 2027 \citep{euAiAct2024}. L'applicabilité à un système donné dépend de son usage et de son rôle juridique.
    \item \textbf{NIST AI RMF :} le référentiel AI RMF 1.0 fournit une base de gouvernance, de cartographie, de mesure et de gestion des risques ; la page officielle indique qu'une révision est en cours \citep{nistAirmf2023}.
\end{itemize}

De cette lecture, je retiens surtout qu'un passage en production demanderait une validation des fonctions risque, conformité, juridique et IT, en complément du travail technique. Les textes seraient à réexaminer à chaque changement de périmètre, de modèle, de fournisseur ou de régime réglementaire.

\section{Grille d'évaluation et scoring}

La grille documentaire utilise sept critères : couverture fonctionnelle, reproductibilité, explicabilité, contrôle du risque, qualité logicielle, intégrabilité et gouvernance. Chaque critère reçoit de 0 à 3 points : absent ou non vérifiable ; partiel ; documenté mais demandant une adaptation ; explicite, testable et directement intégrable. Le maximum est de 21 points. Ces notes rendent mes arbitrages visibles et ne valent pas certification de production.

\begin{table}[H]
\centering
\caption{Scoring documentaire provisoire des dépôts retenus.}
\label{tab:scoring-depots}
\scriptsize
\begin{adjustbox}{width=\textwidth}
\begin{tabular}{lrrrrrrrr}
\toprule
\textbf{Dépôt} & \textbf{Couverture} & \textbf{Reprod.} & \textbf{Explic.} & \textbf{Risque} & \textbf{Qualité} & \textbf{Intégr.} & \textbf{Gouv.} & \textbf{Total} \\
\midrule
LangGraph & 2 & 3 & 2 & 1 & 3 & 3 & 2 & 16/21 \\
Qlib & 3 & 3 & 1 & 1 & 3 & 3 & 1 & 15/21 \\
Lean & 3 & 3 & 1 & 2 & 3 & 3 & 1 & 16/21 \\
OpenBB & 2 & 2 & 1 & 0 & 3 & 3 & 1 & 12/21 \\
TradingAgents & 2 & 1 & 2 & 0 & 2 & 2 & 1 & 10/21 \\
FinRobot & 2 & 1 & 2 & 0 & 2 & 2 & 1 & 10/21 \\
FinRL & 2 & 2 & 1 & 1 & 2 & 2 & 1 & 11/21 \\
Freqtrade & 3 & 2 & 1 & 1 & 3 & 3 & 1 & 14/21 \\
\bottomrule
\end{tabular}
\end{adjustbox}
\end{table}

Le tableau confirme le choix d'une architecture composée. LangGraph obtient son meilleur niveau sur l'orchestration et l'intégrabilité, tandis que Qlib et Lean sont plus solides pour la recherche et le backtesting. Les scores de gouvernance et de contrôle du risque restent faibles ou partiels, ce qui justifie l'ajout de règles métier externes et d'une validation humaine.

\section{Principes de conception}

Ces principes ne sont pas des hypothèses statistiques. Ils fixent les règles de conception qui relient le socle quantitatif, la couche agentique et la gouvernance :

\begin{description}
    \item[Séparation du calcul et de l'interprétation :] les rendements, coûts et indicateurs de risque sont calculés par des fonctions déterministes. Dans la synthèse, les pondérations sont également fixées par le moteur. Dans l'extension d'allocation, le modèle propose des poids ; le code les valide puis calcule les positions et leur PnL.
    \item[Supervision humaine :] une recommandation agentique est un objet de revue. Toute action engageant un portefeuille réel, et notamment l'envoi d'un ordre, reste soumise à une validation explicite et à des permissions indépendantes.
    \item[Explicabilité vérifiable :] une justification doit renvoyer aux données, aux versions, aux contraintes et aux règles réellement utilisées. Le texte produit ne remplace pas le journal structuré.
\end{description}

Pour le classement, j'ai distingué orchestration, recherche financière, infrastructure quantitative et gouvernance. LangGraph et un moteur tel que Lean répondent ainsi à des besoins complémentaires, et non à une même fonction à départager.

\section{Protocole expérimental initial}

J'ai retenu un prototype quantitatif minimal afin de tester le socle déterministe de l'architecture proposée. Cette simplicité est volontaire : l'expérience porte sur la reproductibilité des données, des calculs de portefeuille, des coûts et des indicateurs de risque, plutôt que sur la recherche d'une stratégie optimale.

Le protocole utilise quatre ETF factoriels américains comme proxies investissables de sleeves de valeur, momentum, qualité et faible volatilité : \texttt{VLUE}, \texttt{MTUM}, \texttt{QUAL} et \texttt{USMV}. Ces ETF ne constituent pas une reconstruction titre par titre des facteurs value, momentum, quality et low volatility : ils servent de portefeuilles observables représentant ces expositions. Le benchmark est \texttt{SPY}. Les cours ajustés sont téléchargés via \texttt{yfinance} \citep{yfinance} sur la période allant du 1er janvier 2015 au 31 décembre 2025. Le portefeuille multifactoriel est équipondéré et rééquilibré à la fin de chaque mois. Les prix de fin de mois représentent un prix d'exécution théorique ; aucune exécution au jour suivant, bid--ask spread ou impact de marché n'est simulé. Le portefeuille est supposé déjà investi aux poids cibles à la première date, de sorte que le coût initial n'est pas inclus.

Trois hypothèses de coûts de transaction sont testées : 0, 10 et 30 points de base par unité de turnover. Le taux sans risque est fixé à 0\,\% pour faciliter la comparaison du ratio de Sharpe. Les indicateurs calculés sont le rendement annualisé, la volatilité annualisée, le ratio de Sharpe, le maximum drawdown, la valeur finale d'un investissement initial de 100 unités et le turnover.

Pour un mois $t$, le rendement net du portefeuille est calculé selon :
\[
r_t^{\mathrm{net}} = \sum_i w_{i,t-1}r_{i,t} - c\tau_t,
\qquad
\tau_t = \frac{1}{2}\sum_i\left|w^{\mathrm{cible}}_{i,t}-w^{\mathrm{après rendement}}_{i,t}\right|,
\]
où $c$ désigne le coût unitaire exprimé en fraction et $\tau_t$ le turnover au rebalancement. Sur $T$ rendements mensuels, les métriques principales sont :
\[
R_{\mathrm{ann}} = \left(\prod_{t=1}^{T}(1+r_t)\right)^{12/T}-1,
\qquad
\sigma_{\mathrm{ann}} = \sqrt{12}\,\mathrm{sd}(r_t),
\qquad
S = \sqrt{12}\,\frac{\mathrm{mean}(r_t)}{\mathrm{sd}(r_t)}.
\]
Le turnover est défini en convention « one-way » : il représente la valeur des achats (égale à celle des ventes dans un rééquilibrage neutre), et les coûts sont calculés comme $c\tau_t$. Le turnover annualisé est calculé par $12\sum_t\tau_t/T$. Le drawdown est obtenu à partir du rapport entre la valeur courante et le maximum historique de la valeur cumulée. Ces formules rendent explicite le traitement des coûts, du turnover et de l'annualisation.

Le déroulement peut être résumé par l'algorithme~\ref{alg:backtest}.

\begin{algorithm}[H]
\caption{Calcul du portefeuille multifactoriel et des métriques}
\label{alg:backtest}
\small
\begin{algorithmic}[1]
\Require Cours ajustés des quatre ETF factoriels et de \texttt{SPY}, période d'étude, coûts $\mathcal{C}=\{0,10,30\}$ pb
\Ensure Rapports de qualité, métriques, diagnostics de rééquilibrage et figures
\State Télécharger les cours et vérifier les dates, valeurs manquantes et prix non positifs
\State Supprimer les dates incomplètes et enregistrer \fichiergithub{data/quality_report.json}{quality_report.json}
\State Calculer les rendements mensuels et initialiser les poids cibles à $25\,\%$ par sleeve
\For{chaque coût $c \in \mathcal{C}$}
    \State Réinitialiser le portefeuille aux poids cibles
    \For{chaque mois $t$}
        \State Calculer le rendement brut et les poids après rendement
        \State $\tau_t \gets \frac{1}{2}\sum_i \left|w^{\mathrm{cible}}_{i,t}-w^{\mathrm{après}}_{i,t}\right|$
        \State $r^{\mathrm{net}}_t \gets r^{\mathrm{brut}}_t-c\tau_t$
        \State Enregistrer rendement, turnover, coût et poids
    \EndFor
    \State Calculer rendement annualisé, volatilité, Sharpe, drawdown et valeur finale
    \State Exporter les diagnostics de rééquilibrage pour le coût $c$
\EndFor
\State Comparer le benchmark et le multifactoriel sur l'échantillon complet
\State Répéter la comparaison sur la période hors échantillon 2021--2025
\State Exporter \texttt{metrics.csv}, \texttt{period\_metrics.csv} et les figures
\end{algorithmic}
\end{algorithm}

Le calcul sépare l'échantillon complet de la période hors échantillon allant de janvier 2021 à décembre 2025. Comme aucun paramètre n'est appris sur les rendements, cette coupure constitue un contrôle de robustesse temporelle et non une validation complète d'un modèle prédictif. Les métriques des deux périodes sont conservées dans \texttt{period\_metrics.csv}.

Le déroulement du protocole est résumé par la figure~\ref{fig:protocole-experimental}. Le téléchargement contient 2\,766 observations quotidiennes. Le contrôle qualité n'a relevé ni date dupliquée, ni valeur manquante, ni prix non positif. L'écart calendaire maximal est de quatre jours, aucun gap supérieur à sept jours n'est observé et aucun mouvement quotidien supérieur à 25\,\% n'est signalé. Les contrôles et leurs limites connues sont archivés dans \fichiergithub{data/quality_report.json}{data/quality_report.json}. Les cours sont auto-ajustés par la source, mais les corporate actions, le survivorship bias, la liquidité et le slippage nécessitent encore une validation institutionnelle ; les scénarios de 10 et 30 points de base restent des proxys de coûts tout compris.

\begin{figure}[H]
\centering
\includegraphics[width=\textwidth]{figures/schemas/protocole_experimental.pdf}
\caption{Déroulement du protocole expérimental et contrôle hors échantillon.}
\label{fig:protocole-experimental}
\end{figure}

\section{Protocole d'évaluation de la couche LLM}

L'évaluation de la couche LLM complète le benchmark quantitatif. Elle vérifie qu'un agent peut lire un état quantitatif déjà calculé, restituer ses champs structurés sans les modifier et respecter les permissions et les alertes produites par le workflow déterministe.

Le protocole est exécuté par \texttt{evaluate\_deepseek.py}. Pour chaque cas, le script recharge les mêmes artefacts que le système étudié : \fichiergithub{data/quality_report.json}{data/quality_report.json}, \url{results/monthly_asset_returns.csv} et \url{results/rebalance_diagnostics_10bps.csv}. Il reconstruit ensuite l'état à une date donnée en réutilisant les fonctions de chargement, de contrôle qualité, d'analyse du rééquilibrage et de contrôle du risque. L'évaluation porte sur douze cas : dix dates historiques, un scénario synthétique de turnover élevé entraînant une alerte de risque et un cas où le contrôle qualité bloque explicitement la suite. Dans ce dernier cas, aucun appel API n'est effectué.

Le modèle reçoit exclusivement un objet JSON structuré contenant la date, l'état qualité, les rendements factoriels, le rendement net, le turnover, le coût, les poids cibles, les alertes de risque, les permissions et les contraintes. Il ne reçoit ni les cours bruts, ni un outil de recherche, ni une capacité de placement d'ordre. La sortie attendue est également contrainte par un contrat JSON : résumé, valeurs citées, recommandation descriptive, limites, indicateurs de permission et liste des champs sources cités. Je peux ainsi distinguer la génération linguistique de la décision quantitative.

L'expérience archivée a été réalisée le 7 septembre 2026. Les métadonnées des 33 réponses identifient le modèle \texttt{deepseek-v4-flash} et l'API \url{https://api.deepseek.com}. La configuration enregistrée comporte dix cas historiques, un stress de turnover, un cas bloqué et trois répétitions par cas non bloqué. L'appel défini dans le code utilise une température nulle, une réponse en objet JSON et un plafond de 3\,000 jetons de sortie. Ces paramètres d'appel ne sont pas tous enregistrés dans les métadonnées historiques, ce qui limite la reconstitution exacte de la requête.

Le validateur compare les champs numériques aux valeurs sources avec une tolérance absolue de $10^{-4}$, vérifie la présence des champs attendus et contrôle les indicateurs booléens de permission. L'évaluateur ajoute des contrôles sur les champs supplémentaires, les noms ou préfixes des champs cités, ainsi que la présence de marqueurs d'alertes. Ces vérifications supplémentaires ne sont pas toutes appliquées par le workflow principal. Les contrôles textuels reposent sur des mots-clés ou sur la mention du champ d'alertes ; ils ne vérifient pas la cohérence sémantique de chaque phrase ni la restitution de chaque alerte.

La stabilité des sorties est mesurée sur les valeurs numériques normalisées à quatre décimales et sur les indicateurs de permission ; une mesure distincte compare l'ensemble des champs cités. Le code d'évaluation réessaie les appels qui lèvent une exception, dans une limite de trois tentatives. Une réponse reçue mais rejetée par le score est comptabilisée comme un échec ; le mécanisme actuel ne la régénère pas automatiquement. Les métadonnées archivées ne fournissent pas un journal exhaustif des éventuelles tentatives intermédiaires.

Les résultats sont enregistrés dans \fichiergithub{results/deepseek_evaluation.json}{results/deepseek_evaluation.json} et \texttt{results/deepseek\_evaluation.csv}. Les métadonnées comprennent le modèle, l'URL de l'API, l'identifiant et l'horodatage de requête, la latence et les nombres de jetons. L'archive du 7 septembre ne contient ni empreinte du prompt, ni empreinte des données, ni commit Git. Le code actuel prévoit ces éléments pour de nouvelles traces du workflow DeepSeek, mais cette extension ne complète pas rétroactivement les preuves de l'expérience. La clé d'API est lue depuis \texttt{.env} et n'apparaît pas dans les sorties. Le protocole mesure la conformité aux contrôles disponibles ; il n'évalue pas la performance financière ni la productivité humaine.

\subsection{Analyse complémentaire des notes et du validateur}

Une analyse hors réseau complète l'expérience archivée : elle compare la note déterministe de la dernière date commune avec la première réponse du cas historique correspondant, puis examine les deux autres répétitions. La grille porte sur l'exactitude, la clarté, les alertes, la traçabilité et la longueur ; l'appréciation de lisibilité est qualitative et ne repose pas sur un panel indépendant. Trois altérations portant chacune sur un seul champ et deux témoins de détection sont ensuite appliqués à la première réponse. Ces essais construits évaluent les limites du validateur, sans être assimilés à des erreurs spontanées du modèle. Le script \path{analyze_note_comparison.py} conserve les entrées, les transformations, les scores et les empreintes des sources, sans appel API ni modification de l'évaluation initiale.

\section{Extension : décisions d'allocation et PnL simulé}\label{sec:protocole-allocation}

L'expérience de synthèse ne peut pas modifier le PnL puisqu'elle reprend une allocation déjà calculée. Pour étudier l'effet d'une décision du modèle, j'ai ajouté une seconde expérience : DeepSeek propose des poids cibles, puis un moteur Python indépendant les accepte ou les rejette et calcule les positions. Cette extension utilise le module \path{agent_allocation_backtest.py}. Elle n'exécute pas les sept agents de l'architecture cible et ne passe pas par le graphe LangGraph de restitution. Elle teste un rôle d'allocation LLM encadré par des règles déterministes, sans connexion à un courtier.

\subsection{Période, références et informations disponibles}

La collecte a été réalisée le 12 septembre 2026 sur 56 décisions mensuelles, pour les mois de mai 2021 à décembre 2025. Chaque portefeuille débute avec 100 unités, déjà investies à 25\,\% dans chaque ETF à la clôture du 30 avril 2021. Les anciennes positions supportent la variation jusqu'à la première exécution. Les coûts d'entrée initiale et de liquidation finale sont exclus.

Trois politiques utilisent les mêmes cours ajustés, le même calendrier, les mêmes contraintes et le même moteur de frais : une cible équipondérée de 25\,\% ; une allocation proportionnelle à l'inverse de la volatilité historique, plafonnée à 40\,\% ; et les poids proposés par DeepSeek. Pour la règle inverse volatilité, le poids dépassant le plafond est fixé à 40\,\%, puis le solde est réparti proportionnellement entre les actifs encore libres. La procédure est répétée si nécessaire. Les trois politiques restent soumises au contrôle de turnover, sans ajustement discrétionnaire après un rejet.

À la dernière séance du mois précédent, les entrées comprennent les rendements des quatre ETF sur 21, 63, 126 et 252 séances, la volatilité des 252 derniers rendements quotidiens annualisée avec $\sqrt{252}$, le drawdown sur les 253 derniers cours, les poids courants et les contraintes. Ces horizons sont des approximations en séances de un, trois, six et douze mois. Les prix sont coupés à la date d'observation avant le calcul des indicateurs. Le modèle ne reçoit ni cours d'exécution futur, ni rendement du mois à venir, ni actualité. Un test modifie les cours futurs et vérifie que les entrées demeurent identiques. Cette propriété du programme ne résout pas le biais de préentraînement discuté en section~\ref{sec:biais-allocation}.

\subsection{Contrat, contraintes et traitement d'un rejet}

La sortie doit contenir exactement quatre champs : \path{decision_date}, \path{target_weights}, \path{reason} et \path{cited_fields}. La date doit correspondre à celle des entrées. Les quatre ETF sont obligatoires ; les poids doivent être des nombres finis, non négatifs, dont la somme vaut 1, sans levier ni poche de liquidités. La concentration maximale est de 40\,\% par poids cible. Une tolérance de $10^{-9}$ est admise sur la somme ; la normalisation dans cette seule tolérance est enregistrée. Les poids effectivement appliqués figurent dans le journal.

Le turnover proposé est calculé par $\frac{1}{2}\sum_i|w_i^*-w_i^{\mathrm{courant}}|$ et limité à 10\,\%. Les chemins cités doivent exister exactement dans l'objet transmis, et la justification doit être un texte non vide. Ce validateur d'allocation est distinct de celui des 33 synthèses : il contrôle les contraintes numériques et les chemins complets, mais pas la pertinence économique du raisonnement ni la véracité de chaque phrase.

Une proposition rejetée conserve les \emph{quantités} déjà détenues, sans facturer de transaction. Les poids continuent alors à varier avec les prix ; ils ne sont pas réinitialisés à 25\,\%. Le moteur ne corrige pas les poids hors limites et ne demande pas une autre réponse pour remplacer le rejet. Le plafond de concentration s'applique aux poids cibles, sans garantir que les mouvements de marché le respectent entre deux rééquilibrages. Le choix de conserver les positions fait partie de la politique évaluée.

\subsection{Exécution différée et frais autofinancés}

La proposition est produite à partir de la clôture du mois précédent ; l'exécution simulée intervient à la première clôture disponible du mois suivant. Le turnover est revérifié à ce cours, avant puis après prise en compte des frais. Un dépassement provoque également la conservation des positions. Une proposition admissible à l'observation peut ainsi être inexécutable sous les contraintes retenues.

Soient $q_i$ les quantités avant transaction, $P_i$ les cours ajustés d'exécution, $N_i=q_iP_i$ et $V=\sum_i N_i$. Avec les poids cibles $w_i^*$ et le taux $c=0{,}001$, soit 10 points de base, les frais $F$ satisfont :
\begin{equation}
F=\frac{c}{2}\sum_i\left|w_i^*(V-F)-N_i\right|,
\qquad q_i'=\frac{w_i^*(V-F)}{P_i}.
\label{eq:frais-allocation}
\end{equation}
Le moteur résout cette équation par dichotomie. Le turnover réellement échangé vaut $\frac{1}{2V}\sum_i|q_i'P_i-N_i|$. La valeur des nouvelles positions et les frais doivent vérifier $\sum_i q_i'P_i+F=V$. Les frais sont donc financés par le portefeuille et portent sur la moitié du volume effectivement échangé. Avec ces frais, les montants achetés et vendus ne sont plus exactement égaux.

Les unités fractionnaires obtenues à partir de cours ajustés représentent une comptabilité de rendement total. Elles ne constituent pas des quantités directement utilisables auprès d'un courtier. Aucune simulation séparée du spread, de l'impact de marché, de la liquidité ou de la fiscalité n'est réalisée. Le taux de 10 points de base est une hypothèse expérimentale, sans calibrage empirique sur des transactions réelles.

Les valeurs sont suivies chaque séance. Pour $D$ jours calendaires entre le capital initial et la dernière valorisation, le rendement annualisé est $(V_T/V_0)^{365{,}25/D}-1$. La volatilité et le Sharpe utilisent les rendements quotidiens avec facteur $\sqrt{252}$ et taux sans risque nul ; le drawdown inclut la valeur initiale. Ces conventions diffèrent du backtest initial : période plus courte, exécution différée, frais autofinancés et mesures de risque quotidiennes. Les chiffres des deux expériences ne doivent donc pas être comparés comme s'il s'agissait d'un même protocole. Aucun nouveau scénario LLM à 0 ou 30 points de base n'a été collecté ; ces taux servent uniquement aux tests comptables de cette extension.

\subsection{Configuration du modèle et budget de collecte}

Le modèle demandé est \texttt{deepseek-v4-flash} ; les 56 réponses retournent l'identifiant \texttt{deepseek-flash}. Les deux noms sont archivés, sans possibilité de vérifier la version exacte des poids internes du fournisseur. La campagne utilise une température nulle, le mode JSON, un plafond de 3\,000 jetons de sortie et le réglage explicite \texttt{thinking: disabled}, documenté par DeepSeek \citep{deepseekThinking2026}. La température ne suffit pas à garantir des réponses identiques lors de futurs appels.

La période avait d'abord été envisagée sur les 60 mois de 2021--2025. Quatre tentatives techniques ont précédé la collecte : trois réponses sans JSON ont atteint le plafond de 3\,000 jetons de raisonnement, puis une quatrième requête a été interrompue. Après désactivation de ce mode, les 56 appels restants ont été consacrés à une période commune de mai 2021 à décembre 2025. Ce changement répond à la limite de 60 tentatives et a été décidé avant examen des performances de la campagne corrigée. Les quatre tentatives initiales sont conservées séparément et exclues de ses décisions ; elles restent incluses dans le bilan des ressources consommées.

Il n'y a qu'une réponse par date, sans sélection entre répétitions, ni nouvelle génération après rejet. Chaque tentative est écrite sur disque avant l'appel HTTP et consomme une place du budget, même en cas d'erreur ou d'interruption. Une erreur API arrête la collecte ; une reprise conserve cette tentative comme une absence de proposition. Le protocole ne mesure donc pas la variabilité du PnL entre plusieurs générations.

\subsection{Journalisation et reproduction}

L'archive contient le prompt, les indicateurs, les poids courants, les paramètres, leurs empreintes SHA-256, la réponse brute, les identifiants de requête et de modèle, le statut de fin, les jetons rapportés et la durée. Le journal d'exécution ajoute les quantités avant et après décision, les poids appliqués, les frais, le turnover et les motifs de rejet. Les empreintes du fichier de prix, du code et de l'archive sont également enregistrées.

Le mode de reproduction hors réseau vérifie la correspondance exacte de la requête avant de réutiliser la réponse. Modifier les frais, les poids courants ou le prompt invalide cette correspondance. Il s'agit d'une reproductibilité des calculs à décisions archivées, pas d'une garantie de stabilité du service LLM. Les commandes et fichiers nécessaires figurent en annexe~\ref{ann:allocation}. Un test technique sans LLM utilise une règle de momentum sur 60 mois ; ses performances ne sont pas assimilées à celles de DeepSeek.


\section{Limites de la méthode}

Le benchmark initial évalue un socle quantitatif simple et sa synthèse contrôlée. L'extension mesure le PnL de propositions d'allocation validées, sur une seule trajectoire rétrospective. Aucun des deux dispositifs ne valide un système agentique complet. Un univers plus large, des données historiques point-in-time et une modélisation de la liquidité seraient nécessaires pour étendre la portée financière des résultats.

Les fonctionnalités, licences et conditions de maintenance des dépôts GitHub peuvent évoluer. Leur adoption demande une revue technique datée, une analyse de sécurité des dépendances et une validation juridique de la licence.

% Texte d'exemple du template conservé hors compilation.
